% =============================================================
%  CAPÍTULO I — MARCO CONTEXTUAL DEL PROYECTO TECNOLÓGICO
% =============================================================

\section{Problema de investigación}

\subsection{Planteamiento del problema}

Es innegable la creciente importancia de implementar evaluaciones por competencias curriculares en el contexto educativo. En un mundo cada vez más orientado hacia el desarrollo de habilidades, estas evaluaciones se posicionan como elementos fundamentales para medir con precisión las destrezas de los estudiantes en relación con los objetivos y contenidos del plan de estudios.

Las evaluaciones por competencias se presentan como herramientas esenciales que van más allá de evaluar conocimientos teóricos, abordando la capacidad práctica y la aplicación efectiva de conceptos. Este enfoque integral no solo proporciona una evaluación más precisa del rendimiento estudiantil, sino que también contribuye a preparar a los estudiantes con las habilidades necesarias para enfrentar los retos del entorno educativo actual y futuro.

La implementación efectiva y generalizada de las evaluaciones por competencia presenta desafíos significativos. Aunque la tecnología ha avanzado, la integración de estas evaluaciones en las prácticas educativas sigue siendo limitada. La falta de plataformas especializadas y adaptadas a las necesidades específicas de la evaluación por competencias dificulta su aplicación amplia y eficiente. Asimismo, la diversidad de áreas curriculares y la variabilidad en los estilos de aprendizaje de los estudiantes plantean la necesidad de una solución integral que pueda adaptarse a distintos contextos educativos. La escasez de herramientas que permitan la aplicación de evaluaciones por competencias de manera personalizada y detallada, así como la falta de recursos que faciliten la interpretación y el análisis de los resultados por parte de los docentes, son obstáculos adicionales.

\subsubsection*{Diagnóstico}

Actualmente, el proceso de evaluación en la educación básica elemental, específicamente en el área de matemáticas, se ve limitado por métodos tradicionales que priorizan la memorización sobre el desempeño. Se observa que los docentes carecen de herramientas tecnológicas diseñadas específicamente para medir competencias curriculares de forma individualizada. Esta situación genera una brecha entre lo que el estudiante ``sabe'' teóricamente y lo que puede ``hacer'' en contextos prácticos. Además, la falta de plataformas que centralicen y analicen los resultados de manera detallada impide que el profesorado realice retroalimentaciones oportunas, dejando a los estudiantes con vacíos de aprendizaje que no son detectados a tiempo.

\subsubsection*{Pronóstico}

De no implementarse soluciones tecnológicas que faciliten la evaluación por competencias, persistirá la desconexión entre el currículo académico y las habilidades reales de los estudiantes. Esto resultará en un bajo rendimiento en pruebas estandarizadas y una desmotivación creciente en el aprendizaje de las matemáticas. A largo plazo, la falta de una medición precisa de las destrezas individuales fomentará la exclusión educativa, ya que el sistema no será capaz de adaptarse a los diferentes ritmos y estilos de aprendizaje, limitando la preparación de los estudiantes para los retos académicos y profesionales del futuro.

\subsection{Formulación del problema}

¿Cómo superar los desafíos relacionados con la evaluación de competencias curriculares en matemáticas, asegurando precisión, inclusión y adaptabilidad a las necesidades individuales de los estudiantes de educación básica elemental?

\subsection{Sistematización del problema}

\begin{itemize}[leftmargin=*]
  \item ¿Qué necesidades individuales y criterios deben considerarse en el diseño de una herramienta que permita una medición precisa y eficaz de las competencias de los estudiantes?
  \item ¿Qué elementos utilizar en la formulación de las preguntas para el proceso evaluativo que asegure una evaluación integral de las competencias de los estudiantes?
  \item ¿Qué factores deben considerarse en los procesos de evaluación por competencias para garantizar su efectividad en contextos educativos reales?
\end{itemize}

\section{Objetivos}

\subsection{Objetivo General}

Desarrollar una aplicación web para superar los desafíos relacionados con la evaluación de competencias curriculares en matemáticas, asegurando precisión, inclusión y adaptabilidad a las necesidades individuales de los estudiantes de educación básica elemental aplicando el uso de pictogramas.

\subsection{Objetivos Específicos}

\begin{itemize}[leftmargin=*]
  \item Identificar las necesidades de los estudiantes y las limitaciones que enfrentan los docentes en la evaluación por competencias, mediante entrevistas a profesionales educativos, con el fin de establecer los requisitos funcionales y pedagógicos para el diseño de la aplicación.
  \item Diseñar la estructura y funcionalidades de una aplicación web que integre múltiples tipos de preguntas, ajustes de dificultad basados en el currículo, validación de respuestas y elementos de accesibilidad como pictogramas, con el fin de atender las necesidades individuales de los estudiantes en la evaluación por competencias.
  \item Validar la funcionalidad de la aplicación en entornos educativos reales, observando su desempeño durante la administración de evaluaciones por competencias y recopilando percepciones de docentes y estudiantes para orientar mejoras que fortalezcan su aplicación en el contexto escolar.
\end{itemize}

\section{Justificación}

La evaluación por competencias se ha posicionado como un pilar fundamental en la transformación educativa contemporánea, desplazando el interés desde la mera memorización de contenidos hacia la capacidad de movilizar conocimientos en situaciones prácticas. Sin embargo, en el contexto de la educación básica ecuatoriana, existe una brecha significativa entre estos postulados teóricos y la realidad del aula. Las herramientas actuales, como SCALA o CAT, están diseñadas predominantemente para la educación superior, dejando desatendido el nivel básico elemental. Este vacío instrumental impide que las instituciones educativas locales implementen modelos evaluativos modernos, perpetuando prácticas tradicionales que no logran medir con fidelidad el desarrollo de destrezas críticas en asignaturas fundamentales como las matemáticas.

Desde una perspectiva práctica, el desarrollo de este proyecto responde a la necesidad urgente de dotar a los docentes de un instrumento funcional y especializado que automatice y optimice el proceso evaluativo. La utilidad de la aplicación propuesta radica en su capacidad para ofrecer reportes detallados y personalizados sobre el desempeño del estudiante, algo que manualmente resulta inviable para un profesor con numerosos alumnos.

Por esta razón, se plantea la creación de una aplicación web inclusiva destinada a la realización de evaluaciones por competencias curriculares a estudiantes de primaria, centrada en la disciplina de Matemáticas. La propuesta incluye la segmentación de las evaluaciones por módulos, lo que permite adaptar la dificultad de acuerdo con la estructura de los cursos, así como la integración de una variedad de tipos de preguntas para captar el interés y la participación activa del estudiante. Desde la perspectiva del docente, se subraya la generación de informes individuales detallados para que puedan interpretar y aplicar los resultados en el proceso educativo.

Finalmente, la rigurosidad del proceso de desarrollo se garantiza mediante la adopción de la metodología ágil Extreme Programming (XP). A diferencia de los modelos en cascada o metodologías orientadas al hardware, XP es idónea para este contexto educativo debido a su enfoque en la simplicidad, la retroalimentación constante y la entrega de valor iterativa. Dado que las necesidades pedagógicas pueden ser dinámicas, XP permite validar continuamente las funcionalidades --como los tipos de preguntas y la visualización de reportes-- directamente con los usuarios finales, asegurando que el software resultante no sea solo un ejercicio de programación, sino una solución validada que responde fielmente a los problemas reales detectados en la fase de diagnóstico.
